\documentclass[14pt]{extarticle}
\usepackage[margin=0.85in]{geometry}
\usepackage{enumitem}
\usepackage{listings}
\usepackage{xcolor}
\usepackage{framed}
\usepackage{amssymb}

\lstset{
  basicstyle=\ttfamily\small,
  breaklines=true,
  frame=single,
  backgroundcolor=\color{gray!8}
}

\title{Lecture Note Template Preview}
\author{SteelHacks note-template POC}
\date{}

\begin{document}
\maketitle

\section{C syntax parsing}

\begin{itemize}[leftmargin=*]
  \item C is parsed \hrulefill
  \item Variables and functions must be declared \hrulefill
  \item How is this different from Java? \hrulefill
\end{itemize}

\vspace{1.2em}

\section{Getting input}

\subsection{Command-line arguments}
\begin{itemize}[leftmargin=*]
  \item \texttt{argv[0]} is always \hrulefill
  \item Example program: \hrulefill
\end{itemize}

\subsection{Code: read argv}
\begin{lstlisting}
int main(int argc, char *argv[]) {
    // What does argv[1] contain?
    // ________________________________
}
\end{lstlisting}

\vspace{1.2em}

\section{Streams}

\begin{itemize}[leftmargin=*]
  \item stdin / stdout / stderr are \hrulefill
  \item Why must a stream be flushed? \hrulefill
\end{itemize}

\noindent\textbf{Diagram: buffered stream from program to OS}
\begin{framed}
\begin{minipage}[t][0.28\textheight]{\textwidth}
\vspace{0.5em}
\textit{Sketch the buffer, the flush, and where printf writes.}
\end{minipage}
\end{framed}

\vspace{2em}
\begin{center}
\textit{This template may not cover every topic from the source slides.}
\end{center}

\end{document}
